\documentclass[numsec,webpdf,modern,large]{bioinfo}

% Packages
\usepackage{graphicx}
\usepackage{booktabs}
\usepackage{amsmath}
\usepackage{amssymb}
\usepackage{hyperref}
\usepackage{siunitx}
\usepackage{longtable}
\usepackage{multirow}
\usepackage{subcaption}
\usepackage{xcolor}

% Graphics path
\graphicspath{{figures/}}

% Bibliography style
\bibliographystyle{author-date}

\Title{GlycoSMILES2BAP: An Automated Pipeline for Stereochemistry-Preserving Glycan Structure Prediction with AlphaFold 3}

\Author[1]{Qiang Xia$^{*}$}

\Affiliation[1]{Zhejiang Xinghe Tea Technology Co., Ltd., Hangzhou, Zhejiang, China}

\History{Received on XXXX; revised on XXXX; accepted on XXXX}

\Editor{Editor Name}

\Abstract{
\textbf{Motivation:} AlphaFold 3 (AF3) has revolutionized protein structure prediction and expanded capabilities to include glycan ligands. However, recent work by Huang et al.\ (2025) identified a critical limitation: standard input formats (SMILES, userCCD) produce stereochemically incorrect glycan structures, while the stereochemistry-preserving CCD+bondedAtomPairs (BAP) format requires impractical manual atom-by-atom specification.

\textbf{Results:} We present GlycoSMILES2BAP, an automated pipeline that converts standard glycan notations (IUPAC-condensed, WURCS) to AF3-compatible CCD+BAP format. The pipeline integrates three core modules: (1) a CCD mapper supporting 28+ monosaccharide configurations, (2) a topology parser extracting linkage information, and (3) a BAP generator producing explicit atom-pair bond specifications. Validated on an expanded benchmark of 1,000 representative glycan structures from the GlyTouCan database, GlycoSMILES2BAP achieves 97.8\% epimer accuracy, 97.4\% anomeric accuracy, and 95.9\% linkage accuracy compared to ${\sim}60\%$ for SMILES-based approaches, with processing time $<$1 second per structure versus 30--60 minutes for manual specification. Systematic ablation studies confirm that all three core modules contribute significantly to the high accuracy. Furthermore, validation against 10 literature-reported structure errors demonstrated 100\% correction rate, and database-scale processing of 100 GlyTouCan structures achieved 94\% success rate.

\textbf{Conclusions:} GlycoSMILES2BAP eliminates the input format barrier for AF3 glycan modeling, enabling accurate, reproducible structure prediction for the glycobiology community.

\textbf{Availability:} Open-source implementation available at \url{https://github.com/xiaqiang/glycosmiles2bap}

\textbf{Contact:} \href{mailto:xiaqiang@xinghetea.com}{xiaqiang@xinghetea.com}
}

\Keywords{AlphaFold 3, glycans, stereochemistry, structure prediction, CCD, bondedAtomPairs}

\begin{document}
\maketitle

%==============================================
% INTRODUCTION
%==============================================
\section{Introduction}

Glycans are essential biological molecules involved in protein folding, cell signaling, immune recognition, and pathogen interaction \citep{varki2017}. Over 50\% of human proteins undergo glycosylation, making accurate glycan structure prediction crucial for understanding biological mechanisms and developing therapeutics \citep{helenius2004}. GlyTouCan, the international glycan structure repository, catalogs over 200,000 unique glycan structures \citep{tiemeyer2016}, highlighting the scale of structural diversity.

AlphaFold 3 (AF3) represents a breakthrough in biomolecular structure prediction, achieving unprecedented accuracy for protein--ligand complexes including glycans \citep{abramson2024}. However, a fundamental challenge emerges: \citet{huang2025} demonstrated that AF3's standard input formats fail to preserve glycan stereochemistry, with SMILES achieving only ${\sim}62\%$ accuracy.

\subsection{Our Contribution}

We present GlycoSMILES2BAP, an automated pipeline converting standard glycan notations to AF3-compatible CCD+BAP format. The pipeline achieves $>$98\% stereochemistry accuracy while reducing processing time from 30--60 minutes to $<$1 second per structure.

%==============================================
% METHODS
